\documentclass[11pt]{article}
\usepackage[T1]{fontenc}
\usepackage{lmodern}
\usepackage[margin=0.87in]{geometry}
\usepackage{amsmath,amssymb,amsthm,mathtools}
\usepackage{microtype,booktabs,array,xcolor}
\usepackage[colorlinks=true,linkcolor=black,urlcolor=blue]{hyperref}
\usepackage{fancyhdr}
\hypersetup{pdftitle={Komorebi Space: Effective Groupoid Descent},pdfauthor={Komorebi P.},pdfsubject={Definition, reconstruction theorem, proof and scope}}
\definecolor{ink}{HTML}{20354A}
\pagestyle{fancy}\fancyhf{}
\fancyhead[L]{\small\textsc{Komorebi Space}}
\fancyhead[R]{\small v0.13.0}
\fancyfoot[L]{\small Komorebi P.}
\fancyfoot[R]{\small\thepage}
\setlength{\headheight}{15pt}
\setlength{\parindent}{0pt}
\setlength{\parskip}{6pt}
\setlength{\abovedisplayskip}{9pt}
\setlength{\belowdisplayskip}{9pt}
\newcommand{\V}{\mathcal V}
\newcommand{\K}{\mathfrak K}
\newcommand{\G}{\mathcal G}
\newcommand{\DD}{\operatorname{DD}}
\newcommand{\Res}{\operatorname{Res}}
\newcommand{\ev}{\operatorname{ev}}
\newcommand{\id}{\operatorname{id}}
\newcommand{\source}[1]{\par{\footnotesize\textsf{Lean source:} \texttt{#1}}\par}
\begin{document}
\thispagestyle{plain}
{\Large\color{ink}\textbf{Komorebi Space}}\\[4pt]
{\large Coherent views, residuals, and effective groupoid descent}\\[7pt]
{\small Komorebi P.\quad |\quad Mathematical rundown\quad |\quad Version 0.13.0}
\vspace{10pt}

\section*{1. What the space is}
A Komorebi space is a construction relative to a specified system of mathematical views. A view can record a representation, a resolution level, or both. The arrows specify which changes of view are admissible. The framework does not assume that every representation change is invertible, or that a universal view system has been selected.

Let $\V$ be a category of views and let
\[
F:\V^{op}\longrightarrow\mathbf{Set}
\]
assign possible realizations $F(v)$ to each view. An arrow $f:v\to w$ has a restriction map $f^*:F(w)\to F(v)$, with $(1_v)^*=\id$ and $(g\circ f)^*=f^*g^*$. The strict ambient realization space is
\[
\boxed{\K(F)=\Gamma(F)=
\left\{(x_v)_v\in\prod_v F(v)\;\middle|\;
 f^*(x_w)=x_v\text{ for every }f:v\to w\right\}.}
\]
A point is an entire coherent family, not one preferred presentation. Evaluation
\[
\ev_v:\K(F)\to F(v),\qquad\ev_v(x)=x_v
\]
recovers its appearance at one view. The evaluations jointly determine the point. More strongly, every compatible family of maps $T\to F(v)$ factors uniquely through $T\to\K(F)$. This is the usual limit universal property, realized explicitly in the formalization.

The collection of individual appearances and the space of coherent families are different objects. The former organizes viewwise data. The latter imposes all transition equations. A fiber can be nonempty while no coherent global family exists.

\subsection*{The additive residual and extension problem}
When the fibers are abelian groups and restriction is additive, set
\[
E=\prod_v A_v,\qquad R=\prod_{f:v\to w}A_v,\qquad
(\delta x)_f=f^*(x_w)-x_v.
\]
Then $\Gamma(A)=\ker\delta$. For prescribed coherent data $s$ on $U\subseteq\operatorname{Ob}(\V)$, let $p_U$ restrict a raw family to $U$. With any raw extension $\widetilde s$, define
\[
\mathfrak o_U(s)=[\delta\widetilde s]\in R/\delta(\ker p_U),
\qquad H_U=\ker\delta\cap\ker p_U.
\]
The obstruction is independent of the raw extension. The checked criteria are
\[
\mathfrak o_U(s)=0\iff s\text{ has a global extension},
\]
\[
s\text{ has exactly one extension}\iff
\mathfrak o_U(s)=0\ \text{and}\ H_U=\{0\}.
\]
A nonempty extension fiber is a translated copy of $H_U$. Existence and uniqueness are separate conclusions.
\source{ViewSections.lean; SectionObstruction.lean; SectionExtensionFibers.lean}

\clearpage
\section*{2. Retaining symmetries rather than only solutions}
A set of realizations does not record every equivalence or automorphism. The nonabelian layer uses arbitrary groups $G_v$ with restriction homomorphisms
\[
f^*:G_w\to G_v.
\]
No commutativity or invertibility of these homomorphisms is assumed. A coherent transition system, or cocycle, consists of labels $c_f\in G_v$ satisfying
\[
\boxed{c_{1_v}=1,\qquad c_{g\circ f}=c_f\,f^*(c_g).}
\]
A family $a=(a_v)_v\in\G:=\prod_vG_v$ changes local trivializations by
\[
\boxed{(a\cdot c)_f=a_vc_f[f^*(a_w)]^{-1}.}
\]
This is an action. An arrow labelled $a$, followed by one labelled $b$, has label $ba$. Inverses are given by $a^{-1}$.

The symmetry-retaining Komorebi construction is the action groupoid
\[
\K_G=Z^1(\V;G)//\G,
\qquad
\operatorname{Hom}_{\K_G}(c,d)=\{a:a\cdot c=d\}.
\]
The double slash denotes a groupoid, not just a quotient set. Its objects are cocycles, and its arrows are actual changes of trivialization. The orbit set $Z^1(\V;G)/\G$ remembers only isomorphism classes.

\subsection*{Coherence is not global triviality}
A cocycle determines a twisted set-valued system
\[
T_c(f)(x)=c_f f^*(x).
\]
A global section of that system satisfies
\[
x_v=c_f f^*(x_w).
\]
The cocycle law says that the transition rules compose coherently. It does not assert that a coherent choice of $x$ exists.

For example, take one view with a loop $\sigma$ satisfying $\sigma^2=1$, and constant coefficients $H$. Cocycle labels are involutions $p\in H$. Gauge equivalence is conjugacy of involutions, and automorphisms are centralizers. If $p\ne1$, the regular twisted section equation $x=px$ has no solution, even though the transition cocycle is valid.

\subsection*{What the next theorem must preserve}
Local objects must be allowed to differ on overlaps. What is supplied is an isomorphism between them, not literal equality of their labels. Those isomorphisms must themselves satisfy a compatibility law on triple overlaps. The new theorem glues this full datum while retaining morphisms and automorphisms.

In this version the global gauge groupoid is unframed. Earlier results allow a prescribed framing domain, but the new equivalence below is stated for the empty framing domain. This avoids silently asserting a stronger framed extension.
\source{GroupViewSystems.lean; CocycleGaugeGroupoid.lean; CocycleExamples.lean}

\clearpage
\section*{3. The new effective descent theorem}
Let $(U_i)_{i\in I}$ be subsets of view objects, each regarded as a full subcategory. Assume a \textbf{composition cover}:
\[
\boxed{\forall\bigl(v\xrightarrow f w\xrightarrow g z\bigr),\quad
\exists i:\ v,w,z\in U_i.}
\tag{C}
\]
Every composable pair must be visible inside one patch. Identities make this condition imply coverage of individual objects and arrows.

A descent datum consists of a local cocycle $c_i$ on each patch and overlap gauges $a_{ij}(v)$, for $v\in U_i\cap U_j$, satisfying
\[
\boxed{c_j(f)=a_{ij}(v)c_i(f)[f^*a_{ij}(w)]^{-1}}
\tag{D1}
\]
for overlap arrows $f:v\to w$, together with
\[
\boxed{a_{ik}(v)=a_{jk}(v)a_{ij}(v)}
\tag{D2}
\]
on triple overlaps. The order in (D2) is essential. The identity and inverse laws $a_{ii}=1$ and $a_{ji}=a_{ij}^{-1}$ follow from it.

A morphism from $(c_i,a_{ij})$ to $(c'_i,a'_{ij})$ is a family of local gauges $r_i$ satisfying the local gauge equations and
\[
r_j(v)a_{ij}(v)=a'_{ij}(v)r_i(v).
\]
These data form a groupoid $\DD(U;G)$, with composition and inverses proved in Lean.

\subsection*{Theorem}
Under (C), restriction is an equivalence of groupoids:
\[
\boxed{\Res:\K_G\simeq\DD(U;G).}
\]
Restriction sends a global cocycle to its local restrictions with identity overlap gauges. The theorem says:

\begin{enumerate}
\item Every coherent descent datum is isomorphic to the restrictions of a constructed global cocycle.
\item Every compatible local morphism between restrictions comes from exactly one global morphism.
\end{enumerate}
Thus the theorem is stronger than existence alone, and stronger than a bijection of isomorphism classes. It preserves the arrow structure.

\subsection*{The exact Lean construction}
{\small\ttfamily\raggedright
KomorebiSpace.GroupViewSystem.\\
CocycleDescentData.descentEquivalence\par}
Its result is a Lean categorical equivalence between \texttt{F.GaugeObject} on the empty framing domain and \texttt{F.CocycleDescentData U}. The forward functor is the explicit restriction functor. Full faithfulness and essential surjectivity are proved before invoking mathlib's standard equivalence construction.
\source{DescentDataGroupoid.lean; DescentEquivalence.lean}

\clearpage
\section*{4. How gluing is constructed}
Choose, for each view $v$, a reference patch $\kappa(v)$ containing it. No preferred reference patch is assumed. Define
\[
b_i(v)=a_{\kappa(v),i}(v).
\]
The triple-overlap equation gives
\[
b_j(v)=a_{ij}(v)b_i(v).
\tag{1}
\]
Now express the local cocycles in these reference frames:
\[
\widetilde c_i(f)=b_i(v)^{-1}c_i(f)f^*(b_i(w)).
\tag{2}
\]
Substituting (D1) and (1) shows that $\widetilde c_i(f)=\widetilde c_j(f)$ on overlap arrows. The labels now agree literally. The earlier strict gluing theorem produces one global cocycle $c$.

The local identifications are retained:
\[
\boxed{b_i(v)c_f[f^*(b_i(w))]^{-1}=c_i(f),\qquad b_j=a_{ij}b_i.}
\]
Consequently the original overlap gauges are realized, not replaced by unspecified pairwise equivalences.

\subsection*{Uniqueness with the correct qualification}
An effectuation is a global cocycle together with its local comparison gauges. For two effectuations $(c,b)$ and $(c',b')$, define
\[
\boxed{g_v=b_i'(v)^{-1}b_i(v).}
\]
The overlap laws prove that the expression is independent of the patch containing $v$. It is a global gauge from $c$ to $c'$, and
\[
b'_ig=b_i.
\]
Any comparison satisfying this equation must have the displayed value of $g$. Hence it is unique. These comparisons satisfy the identity and composition laws.

This does not eliminate ordinary automorphisms of $c$. It gives uniqueness \emph{among comparisons compatible with the fixed local identifications}. Changing reference patches can change the literal global representative, but not its isomorphism class as a realization of the descent datum.

\subsection*{Why the full groupoid equivalence follows}
Between the restrictions of two global cocycles, overlap gauges are identities. The overlap equation for a morphism reduces to literal agreement of its local gauge values. The previous gauge-descent theorem glues these values uniquely. This proves full faithfulness.

For an arbitrary descent datum, the constructed $c$ and gauges $b_i$ give an isomorphism from its restrictions to that datum. This proves essential surjectivity. Both properties are established for the explicit restriction functor.
\source{EffectiveCocycleDescent.lean; DescentComparison.lean; DescentEquivalence.lean}

\clearpage
\section*{5. A noncommutative test and the boundary}
Take four ordered views $0\le1\le2\le3$, constant coefficient group $S_3$, and four proper patches
\[
U_i=\{0,1,2,3\}\setminus\{i\}.
\]
Every composable triple lies in a patch, although no patch contains the whole system. The checked example has local labels on a common arrow
\[
c_0(1\to2)=(01),\qquad c_3(1\to2)=(12).
\]
They are unequal. The overlap gauges relate them, and their triple product has the required noncommutative order. The general construction gives a global representative and an actual isomorphism to an independently constructed benchmark.

\subsection*{Pairwise equivalence is not enough}
Let $t=(01)$. Suppose three proposed overlap gauges required
\[
b_1=tb_0,\qquad b_2=tb_1,\qquad b_2=tb_0.
\]
The first two imply $b_2=b_0$, whereas the last would force $t=1$. Thus no such comparisons exist. Every proposed gauge is individually invertible, but the triple law fails. This impossibility is also proved in Lean.

\subsection*{Established scope}
\begin{tabular}{@{}p{0.30\linewidth}p{0.66\linewidth}@{}}
\toprule
\textbf{Established} & Effective gluing through overlap gauges, unique compatible comparisons, and an equivalence of the specified cocycle groupoids under composition covers.\\[5pt]
\textbf{Still separate} & A verified Grothendieck topology, arbitrary site descent, stackification, higher categorical coefficients, analytic gauge estimates, and arithmetic correspondence theorems.\\[5pt]
\textbf{Choice dependence} & The view system is input. A reference frame chooses a representative. The theorem gives controlled equivalence, not an absolutely canonical presentation.\\
\bottomrule
\end{tabular}

\subsection*{Proof, verification, and background}
The five new modules preserve all 78 earlier proof-module sources. The delivery receipt records the build and axiom audits. Independent exact $S_3$ tests check 216 chain cocycles under all 81 reference choices per datum. Those tests supplement the universal Lean proofs.

The abstract descent principle is established mathematics. This continuation integrates it with the actual Komorebi view and gauge constructions. Background definitions are in the Stacks Project, \href{https://stacks.math.columbia.edu/tag/02ZC}{tag 02ZC} and \href{https://stacks.math.columbia.edu/tag/0268}{tag 0268}.

\textbf{Conclusion.} Local descriptions need not be literally identical. Specified, coherent changes of representation suffice for global reconstruction under the proved cover hypothesis, with symmetries retained.
\source{EffectiveDescentExamples.lean; docs/EFFECTIVE\_GROUPOID\_DESCENT.md}
\end{document}
